\section{Arithmetic information and resolved conditional laws}
\label{sec:conditional-laws}
There are two different conditioning operations. An arithmetic event fixes
information about prime signs before the Poisson comparison. Conditioning
on a count or another observable then divides by the probability of that
observable. Keeping the two costs separate avoids treating a rare state as
though an absolute error resolved it automatically.


The following table collects the sufficient budgets. Here $I=-\log\PP(C)$,
$c>0$ is fixed, $Y^\sharp=\lfloor e^V\rfloor$ and
$\widehat Y=\lfloor e^{\widehat V}\rfloor$; $\nu$ and $\widehat\nu$ belong
to those cutoffs. The moving-field intensities $\Lambda$ are at least one.
\begin{center}\small
\begin{tabularx}{\textwidth}{@{}lXl@{}}\toprule
Observable & Sufficient budget & Prime field\\\midrule
Scalar, hard & $I+\log^+\lambda\le V-c\nu$ & $\mathcal F_{Y^\sharp}$\\
Scalar, soft & $2I+\log^+\lambda\le\widehat V-c\widehat\nu$ & $\mathcal F_{\widehat Y}$\\
All exact-level counts & $I+\log\Lambda\le V-c\nu$ & $\mathcal F_{Y^\sharp}$\\
All sites, hard & $I+\log\Lambda+\log(1+\Lambda)\le V-c\nu$ & $\mathcal F_{Y^\sharp}$\\
All sites, optimized & $I+\log\Lambda\le\widehat V/2-c\widehat\nu$ & $\mathcal F_{\widehat Y}$\\\bottomrule
\end{tabularx}
\end{center}
Exact marks and signs can be retained in both field rows. The inequalities
have fixed negative margins; equality at the boundary and optimality for
the model are not asserted. A further conditioning on an observed count
has the separate probability cost analysed below.

\subsection{Stable kernels and information cost}
\begin{lemma}[Stable product lift]\label{lem:stable-lift}
Let $W$ take values in a standard Borel space, let $\mathcal F$ be a
sigma-field, and let $\nu$ be a deterministic target law. If
\[
 \EE\,\dTV(\Law(W\mid\mathcal F),\nu)\le\varepsilon,
\]
then for every $\mathcal F$-measurable random element $V$,
\[
 \dTV(\Law(V,W),\Law(V)\otimes\nu)\le\varepsilon.
\]
For $C\in\mathcal F$ with $\PP(C)>0$, the conditional bound is at most
$\varepsilon/\PP(C)$, both for $W$ and for the joint product lift with $V$.
\end{lemma}
\begin{proof}
Disintegrate the actual and product laws over $\mathcal F$. For a test set,
the difference of its conditional sections is bounded by the conditional
total-variation distance. Integrate, or integrate over $C$ and divide by
$\PP(C)$. Countable configuration spaces used here have measurable kernels;
the standard-Borel formulation permits the additional recorded variables.
\end{proof}

For $C_N\in\mathcal F_{Y_N^\sharp}$ write
$I_N^{\rm info}=-\log\PP(C_N)$. The hard scalar estimate becomes
\begin{equation}\label{eq:hard-conditioned}
 \dTV(\Law(Z_{N,L}\mid C_N),\Pois(\lambda_N))
 \ll \min\{1,e^{I_N^{\rm info}}\lambda_N
           (e^{-V_N+o(\nu_N)}+N^{-1/3+\varepsilon})\}.
\end{equation}
For $C_N\in\mathcal F_{\widehat Y_N}$ the soft proof gives instead
\begin{equation}\label{eq:soft-conditioned}
 \dTV(\Law(Z_{N,L}\mid C_N),\Pois(\lambda_N))
 \ll\min\{1,e^{I_N^{\rm info}-(\widehat V_N-\log^+\lambda_N)/2
                         +o(\widehat\nu_N)}
              +e^{I_N^{\rm info}}N^{-1/3+\varepsilon+o(1)}\}.
\end{equation}
Thus the sufficient scalar budgets are
\[
 I_N^{\rm info}+\log^+\lambda_N\le V_N-c\nu_N,
 \qquad
 2I_N^{\rm info}+\log^+\lambda_N\le\widehat V_N-c\widehat\nu_N.
\]
The hard route can be preferable at a larger information cost, but it
requires measurability with respect to its smaller prime field. These
budgets measure probabilities of events, not the number of primes exposed.
For instance, specifying $k$ independent prime signs costs $k\log2$;
a rank-$k$ compatible affine cylinder has the same cost.

At small intensity the hard estimate is relative: under
$I_N^{\rm info}\le V_N-c\nu_N$ and $\lambda_N\to0$,
\begin{equation}\label{eq:conditional-one-hit}
 \PP(Z_{N,L}>0\mid C_N)=\lambda_N(1+o(1)),\qquad
 \PP(Z_{N,L}=1\mid C_N)=\lambda_N(1+o(1)).
\end{equation}
The polynomial terms are negligible since $I_N^{\rm info}=o(\log N)$.
At large intensity, either valid scalar budget gives a conditional Gaussian
limit after centering by $\lambda_N$ and dividing by $\sqrt{\lambda_N}$.
The target Poisson normal-approximation error must still be added to the
arithmetic error; the quantitative and relative forms are in
\cref{supp:sec:conditional-details}.

\subsection{The exact resolution cost}
\begin{lemma}[Sharp conditioning inequality]\label{lem:conditioning-floor}
Let $\mu,\nu$ be probability laws on $E\times F$ with
$\dTV(\mu,\nu)\le\varepsilon$. Let $B\subset E$ have
$p=\nu_E(B)>\varepsilon$ and write $q=\mu_E(B)$. Then $q>0$ and
\begin{equation}\label{eq:sharp-conditioning}
 \dTV(\mu_F(\cdot\mid B),\nu_F(\cdot\mid B))
 \le\frac{\varepsilon}{\max(p,q)}\le\frac{\varepsilon}{p}.
\end{equation}
\end{lemma}
\begin{proof}
The bound $|p-q|\le\varepsilon$ gives positivity. First suppose $q\le p$.
For any measurable $D\subset F$ put $A=B\times D$. Since $\mu(A)\le q$,
\[
 \frac{\mu(A)}q-\frac{\nu(A)}p
 \le\frac{\mu(A)-\nu(A)+p-q}{p}
 =\frac{\nu(B\times D^c)-\mu(B\times D^c)}p
 \le\frac\varepsilon p.
\]
The supremum of the positive differences equals total variation for
probability laws. If $q\ge p$, interchange the two laws.
\end{proof}
The denominator expresses the correct scale: convergence after conditioning
requires $\varepsilon/p\to0$, not merely $p>\varepsilon$.
This lemma applies to an entire future path at once.

Let $\lambda_{N,m}=2^{\vartheta_N-m}$ and
$p_{N,m}(n)=e^{-\lambda_{N,m}}\lambda_{N,m}^n/n!$.
Write $\varepsilon_{N,m,C}^{\rm path}$ for the total-variation error between
the actual complete future threshold path, under $C_N$, and the target
staircase based at $m$. The aggregated proof works equally for an arbitrary
base threshold in a fixed logarithmic band (including $\lambda_{N,m}<1$).
In its nontrivial information range it gives
\begin{equation}\label{eq:path-error}
 \varepsilon_{N,m,C}^{\rm path}
 \ll e^{I_N^{\rm info}}\lambda_{N,m}
       (1+\log^+(2\lambda_{N,m}))e^{-V_N+o(\nu_N)}
       +N^{-1/3+\varepsilon},
\end{equation}
with an arbitrarily small fixed exponent slack absorbing the enlarged
length and the subpolynomial information factor.

\begin{theorem}[Resolved thinning kernels]\label{thm:resolved-kernel}
If $p_{N,m}(n)>\varepsilon_{N,m,C}^{\rm path}$, then
\begin{equation}\label{eq:resolved-kernel}
 \dTV\left(\Law((Z_{N,b_N+m+j})_{j\ge0}\mid C_N,Z_{N,b_N+m}=n),
                     \operatorname{Thin}_n\right)
 \le\frac{\varepsilon_{N,m,C}^{\rm path}}{p_{N,m}(n)},
\end{equation}
where $\operatorname{Thin}_n$ is the path of $n$ independent geometric
lifetimes from \eqref{eq:target-forward}. If a fixed lower segment is also
retained, the same inequality uses the comparison error based at its
lowest level and gives the reverse immigration kernel
\eqref{eq:target-reverse} jointly with the future.
\end{theorem}
\begin{proof}
Apply \cref{lem:conditioning-floor} to the present count and the entire
retained path. The conditional target is exactly the geometric-lifetime
path by independent Poisson splitting. For lower levels, the increments
below the present level are independent Poisson variables. A finite
intensity half-line comparison cannot control the whole infinite reverse
path, whose total target mass diverges.
\end{proof}
For every integer $n\ge1$ and every $\lambda>0$, the local formula is
\[
 p_\lambda(n)=\frac{e^{-\delta_n}}{\sqrt{2\pi n}}
                 e^{-\lambda h(n/\lambda)},\qquad
 0\le\delta_n\le\frac1{12n},\qquad h(v)=v\log v-v+1.
\]
The correction is independent of $\lambda$; its relative error is at most
$1/(12n)$. For $n=0$ use the exact mass $e^{-\lambda}$ instead.
The telescoping proof is in \cref{supp:sec:conditional-details}.
For $n_N\to\infty$, a sufficient condition for the right side of
\eqref{eq:resolved-kernel} to vanish is
\begin{equation}\label{eq:resolved-budget}
 I_N^{\rm info}+\log^+\lambda_{N,m}
 +\log(1+\log^+(2\lambda_{N,m}))
 +\tfrac12\log n_N+\lambda_{N,m}h(n_N/\lambda_{N,m})
 \le V_N-c\nu_N.
\end{equation}
In the central regime $n_N=\lambda_{N,m}+t_N\sqrt{\lambda_{N,m}}+O(1)$,
$|t_N|=o(\lambda_{N,m}^{1/6})$, this reduces, with a fixed margin, to
\[
 I_N^{\rm info}+\tfrac32\log\lambda_{N,m}+\tfrac12t_N^2
 \le V_N-c\nu_N.
\]
The coefficient $3/2$ is a sufficient error budget, not a threshold of the
model. Its derivation, including the common-band check, is in
\cref{supp:sec:conditional-details}.

\subsection{Quenched comparison on geometric scales}
\begin{corollary}[Almost-sure environmental control]\label{cor:quenched}
Let $N_k\asymp q^k$, $q>1$, and suppose
$\log\Lambda_{N_k}\le V_{N_k}-c\nu_{N_k}$.
For each $0<\beta<c$, the conditional total-variation distance, given
$\mathcal F_{Y_{N_k}^\sharp}$, between the complete aggregated staircase
and its deterministic target is at most $e^{-\beta\nu_{N_k}}$ for all
sufficiently large $k$, almost surely.
\end{corollary}
\begin{proof}
The fibrewise proof of \cref{thm:moving-comparison}(ii) gives mean distance
$e^{-c\nu_N+o(\nu_N)}+N^{-1/3+o(1)}$.
Markov's inequality at $e^{-\beta\nu_N}$ gives a summable series along
$N_k\asymp q^k$, since $\nu_{N_k}\asymp\sqrt{k/\log k}$.
Borel--Cantelli concludes; no independence between scales is used.
\end{proof}
The same argument applies to the labelled field in its own smaller
sufficient range, and to scalar kernels uniformly over the $O(\log N)$
lengths of any fixed logarithmic band with a fixed error margin.
A maximal coupling exists separately in each qualified environment and
scale. This does not furnish a canonical coupling over all scales.
